% Part: sets-functions-relations
% Chapter: functions
% Section: isomorphisms

\documentclass[../../../include/open-logic-section]{subfiles}


\begin{document}

\olfileid{sfr}{fun}{iso}
\olsection{هم‌ریختی}

\begin{explain}
یک \emph{هم‌ریختی} تناظری دوسویی است که ساختار مجموعه‌هایی را که به هم
مرتبط می‌کند حفظ می‌کند؛ و ساختار، به رابطه‌های برقرار میان~!!{element}s
آن مجموعه‌ها بستگی دارد. دو مجموعهٔ~$X=\{1,2,3\}$ و~$Y=\{4,5,6\}$ را
در نظر بگیرید. هر دو مجموعه با رابطه‌های جانشینی، کوچک‌تربودن و
بزرگ‌تربودن ساختار یافته‌اند. هم‌ریختی میان آن دو مجموعه !!a{bijection}
است که این ساختارها را حفظ کند. پس تابع~!!a{bijective} چون
$f \colon X \to Y$ هنگامی هم‌ریختی است که برای همهٔ~$i,j\in X$،
$i<j$ اگر و تنها اگر~$f(i)<f(j)$؛ $i>j$ اگر و تنها اگر
$f(i)>f(j)$؛ و~$j$ جانشین~$i$ باشد اگر و تنها اگر~$f(j)$ جانشین
$f(i)$ باشد.
\end{explain}

\begin{defn}[هم‌ریختی]
فرض کنید~$U$ جفت~$\langle X, R\rangle$ و~$V$ جفت
$\langle Y, S\rangle$ باشد، به‌طوری‌که~$X$ و~$Y$ مجموعه و~$R$ و~$S$
به‌ترتیب رابطه‌هایی بر~$X$ و~$Y$ باشند. !!{bijection}~$f$ از~$X$
به~$Y$، هم‌ریختی‌ای از~$U$ به~$V$ است اگر و تنها اگر ساختار رابطه‌ای را
حفظ کند؛ یعنی برای هر~$x_{1},x_{2}\in X$،
$\tuple{x_1,x_2} \in R$ اگر و تنها اگر
$\tuple{f(x_1),f(x_2)} \in S$.
\end{defn}

\begin{ex}
دو مجموعهٔ~$X=\{1,2,3\}$ و~$Y=\{4,5,6\}$ و رابطه‌های کوچک‌تربودن و
بزرگ‌تربودن را در نظر بگیرید. تابع~$f\colon X \to Y$ که
$f(x) = 7-x$، هم‌ریختی‌ای میان~$\tuple{X,<}$ و~$\tuple{Y,>}$ است.
\end{ex}

\end{document}
